\documentclass[11pt]{article}
\usepackage[utf8]{inputenc}
\usepackage{amsmath,amssymb,amsthm,enumitem}
\usepackage[margin=1in]{geometry}
\usepackage{xcolor}
\usepackage{hyperref}
\hypersetup{colorlinks=true,linkcolor=blue!50!black}

\newtheorem{theorem}{Theorem}[section]
\newtheorem{lemma}[theorem]{Lemma}
\newtheorem{proposition}[theorem]{Proposition}
\newtheorem{corollary}[theorem]{Corollary}
\theoremstyle{definition}
\newtheorem{definition}[theorem]{Definition}
\newtheorem{problem}[theorem]{Problem}
\theoremstyle{remark}
\newtheorem{remark}[theorem]{Remark}
\theoremstyle{definition}
\newtheorem{conjecture}[theorem]{Conjecture}
\newtheorem{hypothesis}[theorem]{Hypothesis}

\newcommand{\Z}{\mathbb{Z}}
\newcommand{\Q}{\mathbb{Q}}
\newcommand{\R}{\mathbb{R}}
\newcommand{\C}{\mathbb{C}}
\newcommand{\Isom}{\operatorname{Isom}}
\newcommand{\cd}{\operatorname{cd}}

\title{Growth, collisions and lamplighter structure of\\ generic triangle and orthoscheme reflection groups}
\author{Vico Bonfioli}
\date{\today}

\begin{document}
\maketitle

\begin{abstract}
This paper consolidates a group of preprints on the reflection groups of generic Euclidean
orthoschemes and their growth. The unconditional core is: the right-angled Coxeter envelope
$W_n$ has rational growth $(1+t)^{\lfloor n/2\rfloor+1}/K_{n+1}(t)$ with rate
$r_n=1+2\cos(2\pi/(n+3))$; a dichotomy identifies the collision depth of a leg tuple with half
the shortest kernel length; for $n\ge3$ every length-six kernel element lies in a rank-two
standard parabolic and occurs only at dihedral angle $\pi/3$; over $\Q_{>0}^n$ the rank-two
kernel condition is decided by the positional statistics of the leg tuple, the Diophantine input
being three rank-zero elliptic curves, and the corresponding statement over $\R$ is false. In the
plane the kernel of the triangle representation is nontrivial and the growth series of the
generic triangle group departs from the envelope at depth $11$. The growth rate of the planar
right-triangle group is $\beta_2=1/\sqrt{q^\ast}=1.4916177871\ldots$, where $q^\ast$ is the
smallest positive zero of a Hahn--Exton $q$-cosine; $\beta_2<3/2$ is unconditional. The travel
resolvent is transcendental over $\Q(x)$ unconditionally.

Section~\ref{sec:reduce} contains the one statement here that is not in the source preprints:
for every $n\ge2$, any embedding of $W_n$ into $\Isom(\R^n)$ has injective linear part, so the
ambient embedding problem is equivalent to the question of whether $W_n$ embeds in the
\emph{compact} group $O(n)$. At $n=2$ this recovers the planar obstruction from solvability of
$O(2)$ alone, without amenability of $\Isom(\R^2)$, and it isolates non-commutativity of $O(n)$
as the exact point at which the planar argument fails for $n\ge3$.

Section~\ref{sec:status} states, for each result, what it is conditional on. Three hypotheses are
open and are named wherever they are used: the strand-walk transfer model $(M)$, the junction
pairing $(R\!-\!J)$, and the invariance of the travel block $(T)$. Transcendence of the full
growth series $U$ and of its relaxed companion $V$ is conditional on these and is not claimed
here. Irrationality of $\beta_2$ is open.

Sections~\ref{sec:masterC} and \ref{sec:faithful} settle the three-dimensional case. The
conjecture that the orthoscheme representation is faithful at every Class C rational leg tuple is
false, with explicit counterexamples, the shortest kernel element found having length $76$. On the
other hand the tuple $(1,2,11)$ is faithful, proved by identifying $W_3$ as
$F_2\rtimes(\Z/2\times\Z/2)$ and recognising the image of the free factor as the classical
free group of rotations by $\arccos(3/5)$ about perpendicular axes. Consequently, in dimension
three, $\bigcap_a\ker\rho_a$ is trivial, the generic growth series is the rational envelope in
every degree, $W_3$ embeds in $\Isom(\R^3)$, and the generic $3$-orthoscheme group is finitely
presented.
\end{abstract}

\section{What is new, what is known, and what is conditional}\label{sec:status}

\subsection{Prior art}
The following are not ours and are used as such. Isola and Piergallini proved that the generic
triangle reflection group has free metabelian rotation subgroup of rank two and is not finitely
presented. Droms, Lewin and Servatius, and Myasnikov, Roman'kov, Ushakov and Vershik, gave the
length formula $\|\phi\|_1+2\,\mathrm{st}(\phi)$ for free solvable quotients $F/N'$; Sale gave
the metric behaviour of the Magnus embedding. Steinberg's formula for Poincar\'e series of
Coxeter groups is standard, as are the independence polynomials of paths and the location of
their roots. Adamczewski, Dreyfus and Hardouin proved hypertranscendence of the solutions of the
relevant $q$-difference equation for all $q$ not a root of unity, which is stronger than the
$SL_2$ statement recorded in \S\ref{sec:qcos}; our increment there is only the removal of a
transcendence hypothesis on $q$ in that one special case.

\subsection{Claim labels}
Every numbered statement below carries exactly one of PROVED (complete human proof),
VERIFIED (formalized in Lean 4, axioms audited), HEURISTIC (computational support only), or
CONJECTURE. The effective strength of a result is the weakest label in its dependency chain.

\subsection{The three open hypotheses}
\begin{hypothesis}[transfer model, $(M)$]\label{hyp:M}
The strand-walk model computes the relaxed word length, that is, the crossing optimisation of the
site-cost chain is the relaxed length.
\end{hypothesis}
\begin{hypothesis}[junction pairing, $(R\!-\!J)$]\label{hyp:RJ}
The junction pairing $\Pi_y(q_m)$ is nonzero at infinitely many travel poles $q_m$.
\end{hypothesis}
\begin{hypothesis}[travel-block invariance, $(T)$]\label{hyp:T}
The denominator $1-\Sigma_1$ is the identical, cycle-marker-free object in the generating
functions of $U$ and of $V$, so $U$ and $V$ share the travel poles.
\end{hypothesis}

Hypothesis~\ref{hyp:T} is HEURISTIC. Its strand-walk argument rests on the localisation of
cycles, checked exhaustively on the $3869$ pure-travel elements of true length at most $16$, and
on the metric identity $\ell_T=\ell_R+2c$, whose upper bound is PROVED and whose lower bound is
open. Hypothesis~\ref{hyp:RJ} is CONJECTURE. Hypothesis~\ref{hyp:M} is HEURISTIC and everything
in the block-assembly sections is conditional on it.

\begin{remark}[what is not claimed]\label{rem:notclaimed}
Three claims that appeared in a draft consolidation are withdrawn here because they are not
supported by the sources.
\begin{enumerate}[label=(\roman*),leftmargin=2.2em]
\item Irrationality of $\beta_2$ is \emph{open}. It is not proved, and no effective irrationality
measure for $\beta_2$ is available. What is proved is $\beta_2\ne3/2$, by the squeeze of
\S\ref{sec:beta}. The reformulation known is that $\beta_2$ is irrational if and only if a
Hahn--Exton function omits the value $0$ at the relevant rational points; that is a restatement,
not a proof, and it supplies no exponent.
\item Hypothesis~\ref{hyp:T} is \emph{not} a theorem.
\item The lower bound of the metric identity $\ell_T=\ell_R+2c$ is \emph{open}; only the upper
bound is PROVED. Statements that consume the identity inherit that standing.
\end{enumerate}
\end{remark}

\section{The envelope and the dichotomy}\label{sec:env}

Let $O(a_1,\dots,a_n)\subset\R^n$ be the right-corner orthoscheme with vertices $V_0=0$ and
$V_k=V_{k-1}+a_ke_k$, and let $\widehat R_0,\dots,\widehat R_n$ be the reflections in its facets,
generating $W(a)\le\Isom(\R^n)$. Facets $i,j$ with $|i-j|\ge2$ are perpendicular for every leg
tuple, so with $\Gamma_n$ the graph on $\{0,\dots,n\}$ having an edge whenever $|i-j|\ge2$,
\begin{equation}\label{eq:Wn}
W_n=\bigl\langle R_0,\dots,R_n\ \bigm|\ R_i^2=1,\ (R_iR_j)^2=1\ (|i-j|\ge2)\bigr\rangle
\end{equation}
is the right-angled Coxeter group on $\Gamma_n$, and $\rho_a\colon W_n\twoheadrightarrow W(a)$.

\begin{theorem}[Envelope; PROVED]\label{thm:envelope}
$W_n(t)=(1+t)^{\lfloor n/2\rfloor+1}/K_{n+1}(t)$, with $K_0=K_1=1$,
$K_{2j}=K_{2j-1}-tK_{2j-2}$ and $K_{2j+1}=(1+t)K_{2j}-tK_{2j-1}$. The growth rate is
$r_n=1+2\cos(2\pi/(n+3))$.
\end{theorem}

\begin{theorem}[Dichotomy; PROVED]\label{thm:dichotomy}
For a marked quotient $Q=G/N$ with $K(N)=\min_{1\ne g\in N}\ell_G(g)$ and $h=\lceil K/2\rceil$,
the sphere sizes satisfy $s_d(Q)=s_d(G)$ for $d<h$ and $s_h(Q)<s_h(G)$. Hence
$\cd_n(a)=K(\ker\rho_a)/2$.
\end{theorem}

\begin{theorem}[Length six; PROVED]\label{thm:len6}
For $n\ge3$ and $a\in\R_{>0}^n$, if $w\in\ker\rho_a$ has $\ell_{W_n}(w)=6$ then
$w=(R_iR_{i+1})^{\pm3}$ for some $i$ with $\theta_{i,i+1}=\pi/3$.
\end{theorem}

\begin{theorem}[Rank-two classification; PROVED]\label{thm:rank2}
For $a\in\Q_{>0}^n$, $\ker\rho_a$ meets a rank-two standard parabolic nontrivially if and only if
$e(a)+t(a)\ge1$, where $e$ counts endpoint equal pairs and $t$ counts three consecutive equal
legs. The Diophantine input is the triviality of the rational points on three quartics with
Jacobians $24a1$, $72a2$ and $y^2=x^3-x$, all of rank zero. Over $\R$ the statement is false:
$a=(\sqrt3,1,5)$ has a rank-two kernel element with no equal legs.
\end{theorem}

\section{The ambient embedding reduces to a compact group}\label{sec:reduce}

This section is the one place where the present paper goes beyond the sources. Write
$L\colon\Isom(\R^n)\to O(n)$ for the linear part, with kernel the translation group $\R^n$.

\begin{lemma}[PROVED]\label{lem:noab}
For every $n\ge2$, $W_n$ has no nontrivial normal abelian subgroup.
\end{lemma}

\begin{proof}
The Coxeter diagram of the system \eqref{eq:Wn} has an edge exactly where $m_{ij}\ge3$, that is,
exactly where $\{i,j\}$ is a non-edge of $\Gamma_n$, that is, exactly where $|i-j|=1$. So the
diagram is the path on $n+1$ nodes with every edge labelled $\infty$. The path is connected,
so $(W_n,S)$ is irreducible. Every label is $\infty$, so $W_n$ is infinite and is not of affine
type, affine Coxeter systems having all labels finite.

Let $A\trianglelefteq W_n$ be abelian; $A$ is amenable, so it lies in the amenable radical of
$W_n$. We show that radical is trivial, in two steps.

First it is finite. An irreducible, infinite, non-affine Coxeter group contains a rank-one
isometry for its action on the Davis complex (Caprace--Fujiwara), and a group acting properly
cocompactly on a CAT(0) space with a rank-one isometry, and which is not virtually cyclic, is
acylindrically hyperbolic. $W_n$ is not virtually cyclic: for $n\ge3$ it contains
$\langle R_0,R_2\rangle\times\langle R_1\rangle$, and for $n=2$ it is
$(\Z/2\times\Z/2)\ast\Z/2$, which contains a free group of rank two. For an acylindrically
hyperbolic group the amenable radical coincides with the finite radical, the unique maximal
finite normal subgroup (Osin). Note that acylindrical hyperbolicity gives finiteness here, not
triviality; the second step is needed.

Second it is trivial. The finite radical of a Coxeter group is the direct product of its finite
irreducible components, and $W_n$ is irreducible and infinite, so it has no finite component and
its finite radical is trivial. Hence $A=1$.
\end{proof}

\begin{proposition}[Reduction to $O(n)$; PROVED]\label{prop:reduce}
Let $n\ge2$ and let $\varphi\colon W_n\to\Isom(\R^n)$ be an injective homomorphism. Then
$L\circ\varphi$ is injective. Consequently
\[
  W_n\ \text{embeds in}\ \Isom(\R^n)\quad\Longleftrightarrow\quad W_n\ \text{embeds in}\ O(n).
\]
\end{proposition}

\begin{proof}
$\ker(L\circ\varphi)=\varphi^{-1}\bigl(\varphi(W_n)\cap\R^n\bigr)$ is normal in $W_n$ and is
isomorphic to a subgroup of the abelian group $\R^n$, hence abelian. By Lemma~\ref{lem:noab} it is
trivial. The equivalence follows, one direction being $O(n)\le\Isom(\R^n)$.
\end{proof}

\begin{corollary}[The plane, without amenability; PROVED]\label{cor:plane}
$W_2$ does not embed in $\Isom(\R^2)$.
\end{corollary}

\begin{proof}
$O(2)=SO(2)\rtimes\Z/2$ is metabelian, hence solvable, hence contains no free subgroup of rank
two. $W_2=(\Z/2\times\Z/2)\ast\Z/2$ is a free product of finite groups of orders $4$ and $2$ and
is not infinite dihedral, so it contains a free subgroup of rank two. So $W_2$ does not embed in
$O(2)$, and Proposition~\ref{prop:reduce} transfers this to $\Isom(\R^2)$.
\end{proof}

\begin{remark}[where the planar argument really fails]\label{rem:wherefails}
The published planar obstruction runs through amenability of $\Isom(\R^2)$. That route is
unavailable for $n\ge3$ because $SO(3)$ contains a free subgroup of rank two, so $\Isom(\R^n)$
is not amenable as an abstract group. Proposition~\ref{prop:reduce} relocates the obstruction:
the translation part is never the issue in any dimension, and the planar proof uses exactly one
property of the ambient group, namely that $O(2)$ is solvable. For $n\ge3$ the target $O(n)$ is a
compact, non-solvable Lie group, and that single change is the whole of the difficulty.
\end{remark}

\begin{remark}[what the reduction buys]\label{rem:buys}
Replacing $\Isom(\R^n)$ by $O(n)$ replaces a non-compact target by a compact one, and compactness
is usable. If $W_n\le O(n)$, its closure $H$ is a compact Lie subgroup of $O(n)$, $W_n$ is dense
in $H$, and $H$ is infinite since $W_n$ is. Hence $H^\circ$ is a connected compact Lie group of
dimension at most $n(n-1)/2$, the subgroup $W_n\cap H^\circ$ has finite index in $W_n$ and is
dense in $H^\circ$. In particular an embedding would give $W_n$ a faithful orthogonal
representation of degree $n$. This is a much more rigid demand than a faithful action by
isometries of $\R^n$, and it is the form in which we propose Problem~\ref{prob:embed} be attacked.
\end{remark}

\begin{remark}[the point group, measured]\label{rem:pointgroup}
Proposition~\ref{prop:reduce} says the linear part carries everything, so the object to measure
is the point group $\pi\circ\rho_a\colon W_n\to O(n)$, not $\rho_a$. If the point group is
injective then so is $\rho_a$, so a faithful point group would settle
Problems~\ref{prob:alternative}, \ref{prob:embed} and \ref{prob:fp} together; the converse
fails, and a deviation of the point group is not a deviation of $\rho_a$. The point group has not
previously been measured on its own. Running it against the envelope at two independent primes
$2^{61}-1$ and $2^{62}-57$, at Class C leg tuples, gives agreement to the following depths:
\[
\begin{array}{r|cccccc}
 n & 2 & 3 & 4 & 5 & 6\\\hline
 \text{legs} & (1,2) & (1,2,3) & (1,2,3,5) & (3,1,4,2,5) & (1,2,3,5,7,4)\\
 \text{agrees to depth} & \textbf{2} & 22 & 17 & 15 & 13
\end{array}
\]
and a second tuple at $n=3$, namely $(2,5,3)$, and at $n=4$, namely $(3,1,4,2)$, give the same
depths as the first. The $n=2$ column is a negative control and it fires: there the point group
lies in $O(2)$, which is solvable, while $W_2$ contains a free group of rank two, so a deviation
is forced, and it occurs at depth $3$, the earliest depth not already fixed by the
$(R_0R_2)^2=1$ relation. For $n\ge3$ no deviation is found. This is HEURISTIC and decides
nothing: it is consistent both with faithfulness and with a horizon beyond the computed range.
What it does establish, at each tuple listed, is that the point group is injective on the ball of
the stated radius, hence so is $\rho_a$; at $n=3$ that improves the previously recorded depth for
$\rho_a$ from $19$ to $22$, and does so for the stronger object.
\end{remark}

\section{Class C faithfulness is false}\label{sec:masterC}

Remark~\ref{rem:pointgroup} found Class C tuples at which the point group is \emph{not}
injective. That is not by itself a statement about $\rho_a$, the point group being a quotient.
The following elementary observation converts it into one, and needs neither
Lemma~\ref{lem:noab} nor any CAT(0) input.

\begin{proposition}[Translations commute; PROVED]\label{prop:transtrick}
Let $a\in\R_{>0}^n$ and let $u\in W_n$ satisfy $\pi(\rho_a(u))=I$, that is, $\rho_a(u)$ has
trivial linear part. Then for every $g\in W_n$,
\[
   [\,u,\ gug^{-1}\,]\ \in\ \ker\rho_a .
\]
\end{proposition}

\begin{proof}
An isometry of $\R^n$ with trivial linear part is a translation, so $\rho_a(u)$ is a translation.
For any $g$, $\rho_a(gug^{-1})=\rho_a(g)\rho_a(u)\rho_a(g)^{-1}$ is a conjugate of a translation
by an isometry, hence again a translation. Two translations of $\R^n$ commute, so the commutator
of their preimages is killed.
\end{proof}

\begin{theorem}[Conjecture C is false; PROVED]\label{thm:masterCfalse}
There are Class C tuples $a\in\Q_{>0}^3$, that is with pairwise distinct coordinates, at which
$\rho_a$ is not injective. Explicitly this holds at
\[
 (1,3,2),\quad (2,3,9),\quad (2,6,9),\quad (3,1,5),\quad (3,2,6),\quad (4,9,7),\quad (6,11,7)
\]
and at their reversals. At $a=(3,1,5)$ an explicit nontrivial element of $\ker\rho_a$ has
length $76$, obtained by reducing an explicit word of length $92$.
\end{theorem}

\begin{proof}
Take $a=(3,1,5)$; the other tuples are identical in form. Breadth-first search in $W_3$, using
the right-angled normal form and exact rational arithmetic for the linear reflections, produces
two elements $w_1\ne w_2$ of $W_3$ of length $11$ with the same linear part; the sphere count of
the point group at depth $11$ falls short of the envelope, at each of two independent primes.
Put $u=w_1w_2^{-1}$, of length $22$. Then $\pi(\rho_a(u))=I$ by construction, and $u\ne1$ in
$W_3$, verified in the Tits geometric representation of the Coxeter system, which is faithful for
every Coxeter system and has integer matrices here, so the verification is exact. By
Proposition~\ref{prop:transtrick}, $c=[u,\,R_1uR_1]$ lies in $\ker\rho_a$, and $c\ne1$ in $W_3$,
again by the Tits representation; reducing $c$ gives a geodesic of length $76$. Finally the
coordinates of $a$ are pairwise distinct, so $a$ is Class C.
\end{proof}

\begin{remark}[why this was not seen]\label{rem:whynotseen}
The counterexamples are invisible to a direct search. By the dichotomy the deviation of $\rho_a$
sits at depth $\lceil K/2\rceil$, and here $K\le76$, so the first depth at which the orbit growth
of $W(a)$ departs from the envelope is somewhere below $38$ but far above any reachable
breadth-first search: the envelope sphere at depth $38$ is $9\cdot2^{36}$, already above
$6\cdot10^{11}$. A direct
computation on $\rho_a$ at these tuples sees agreement with the envelope for as long as it can
run, and indeed does. What makes the counterexamples accessible is that the \emph{linear part}
fails much earlier, at depth $11$ to $13$, where the search is trivial. That is the practical
content of Proposition~\ref{prop:reduce}: the obstruction lives in $O(n)$, and it can be seen
there.
\end{remark}

\begin{corollary}[a necessary condition; PROVED]\label{cor:necessary}
If $\rho_a$ is injective then the point group $\pi\circ\rho_a$ is injective. Equivalently, a
point-group collision at $a$ certifies that $\rho_a$ is not injective.
\end{corollary}

\begin{proof}
$\ker(\pi\circ\rho_a)$ is normal in $W_n$, and if $\rho_a$ is injective it is isomorphic to its
image, a subgroup of the translation group, hence abelian. By Lemma~\ref{lem:noab} it is trivial.
\end{proof}

\begin{remark}[what remains open]\label{rem:stillopen}
Theorem~\ref{thm:masterCfalse} settles the arithmetic form of Problem~\ref{prob:alternative}
negatively. It does \emph{not} settle the generic form: the tuples above are single points, and
$N_{\mathrm{gen}}=\bigcap_a\ker\rho_a$ may still be trivial, since a kernel element at one tuple
need not be a kernel element at another. It also leaves Problem~\ref{prob:embed} open, an
abstract embedding not being required to come from an orthoscheme.
\end{remark}

\section{A faithful tuple in dimension three}\label{sec:faithful}

\begin{lemma}[Structure of $W_3$; PROVED]\label{lem:W3struct}
Put $A=R_0R_1$, $C=R_2R_3$ and $V=\langle R_0,R_3\rangle$. Then $V\cong\Z/2\times\Z/2$,
$\langle A,C\rangle$ is free of rank two with basis $\{A,C\}$, and
$W_3=\langle A,C\rangle\rtimes V$, the action being
$R_0AR_0=A^{-1}$, $R_0CR_0=C$, $R_3AR_3=A$, $R_3CR_3=C^{-1}$.
\end{lemma}

\begin{proof}
The four conjugation identities are immediate from the defining relations, using that $R_0$
commutes with $R_2,R_3$ and $R_1$ with $R_3$. Hence $V$ normalises $\langle A,C\rangle$, and since
$R_1=R_0A$ and $R_2=CR_3$ the product $\langle A,C\rangle V$ is all of $W_3$.

Let $\varphi\colon W_3\to(\Z/2)^2$ send $R_0,R_1\mapsto(1,0)$ and $R_2,R_3\mapsto(0,1)$; the
defining relations are respected, so $\varphi$ is well defined, and it is onto. Every nontrivial
element of a finite parabolic $\langle R_i,R_j\rangle$ has nonzero image, and every torsion
element of a right-angled Coxeter group is conjugate into a finite parabolic, so $\ker\varphi$ is
torsion free. As $\Gamma_3$ is the path on four vertices, $W_3$ is the fundamental group of a
graph of groups with three vertex groups $\Z/2\times\Z/2$ and two edge groups $\Z/2$, so it is
virtually free with $\chi(W_3)=3/4-2/2=-1/4$. A torsion-free virtually free group is free, so
$\ker\varphi$ is free, of index $4$ and hence of Euler characteristic $-1$, that is of rank two.
Finally $\langle A,C\rangle\subseteq\ker\varphi$, $\langle A,C\rangle\cap V=1$ because
$\varphi$ is injective on $V$, and $\langle A,C\rangle V=W_3$; so $\langle A,C\rangle$ has index
$4$ and equals $\ker\varphi$. A free group of rank two generated by two elements has them as a
basis.
\end{proof}

\begin{theorem}[A faithful tuple; PROVED]\label{thm:faithful}
At $a=(1,2,11)$ the point group $\pi\circ\rho_a\colon W_3\to O(3)$ is injective. Hence $\rho_a$ is
injective.
\end{theorem}

\begin{proof}
Write $\alpha$ for the angle with $\cos\alpha=3/5$, $\sin\alpha=4/5$. At this tuple
$|m_1|^2=5$ and $|m_2|^2=125$, and a direct computation gives
\[
 \pi(\rho_a(A))=\begin{pmatrix}3/5&-4/5&0\\4/5&3/5&0\\0&0&1\end{pmatrix},\qquad
 \pi(\rho_a(C))=\begin{pmatrix}1&0&0\\0&-117/125&-44/125\\0&44/125&117/125\end{pmatrix},
\]
the first being the rotation about $e_3$ by $\alpha$ and the second the rotation about $e_1$ by
$3\alpha$, since $\cos3\alpha=4\cos^3\alpha-3\cos\alpha=-117/125$ and
$\sin3\alpha=3\sin\alpha-4\sin^3\alpha=44/125$. So with $B$ the rotation about $e_1$ by $\alpha$,
which is not itself in the group, $\pi(\rho_a(C))=B^3$.

The rotations by $\alpha$ about two perpendicular axes generate a free group of rank two. This is
the classical free-rotation theorem in exactly the form stated for $\cos\alpha=3/5$
\textup{(}Wagon, \emph{The Banach--Tarski Paradox}, Thm.~2.1\textup{)}, proved by the $5$-adic
argument: for a reduced word $w$ of
length $n$ whose first applied letter is $A^{\pm1}$, the vector $5^n\,w\,(1,0,0)^{T}$ is integral
and its second coordinate is not divisible by $5$, so $w\ne I$; a word not of that form is
conjugate to one that is, or is a power of $B$, which has infinite order because $3/5$ is not the
cosine of a rational multiple of $\pi$ (Niven). The invariant was checked independently over all
$4{,}782{,}968$ reduced words of length at most $14$ of that form, in exact integer arithmetic. A
subgroup of a free group is free, and $A$ and $B^3$ form a basis of a free group of rank two,
since a reduced word in $A$ and $B^3$ substitutes to a word alternating $A$-blocks with
$B$-blocks of nonzero exponent, which is already reduced and nonempty. Hence $\langle\pi\rho_a(A),\pi\rho_a(C)\rangle$ is
free of rank two on those two elements, so by Lemma~\ref{lem:W3struct} the restriction of
$\pi\circ\rho_a$ to $\langle A,C\rangle$ is injective.

The image of $V$ consists of the four distinct diagonal matrices
$I,\ \mathrm{diag}(-1,1,1),\ \mathrm{diag}(1,1,-1),\ \mathrm{diag}(-1,1,-1)$, so $V$ injects. The
image of $\langle A,C\rangle$ is free, hence torsion free, while every nontrivial element of the
image of $V$ is an involution, so the two images intersect trivially. Now if
$\pi\rho_a(kv)=1$ with $k\in\langle A,C\rangle$ and $v\in V$, then
$\pi\rho_a(v)=\pi\rho_a(k)^{-1}$ lies in that intersection, so $\pi\rho_a(v)=1$ and
$\pi\rho_a(k)=1$, whence $v=1$ and $k=1$. So $\pi\circ\rho_a$ is injective, and
$\ker\rho_a\subseteq\ker(\pi\circ\rho_a)=1$.
\end{proof}

\begin{corollary}[Three consequences in dimension three; PROVED]\label{cor:three}
For $n=3$:
\begin{enumerate}[label=(\roman*),leftmargin=2.2em]
\item $N_{\mathrm{gen}}=\bigcap_a\ker\rho_a=1$, so $\rho_a$ is injective for every $a$ in a
co-null set of leg tuples, and the generic $3$-orthoscheme growth series is the rational envelope
$W_3(t)=(1+t)^2/(1-2t)$ in every degree;
\item $W_3$ embeds in $O(3)$, hence in $\Isom(\R^3)$, answering the ambient embedding question
affirmatively;
\item the generic $3$-orthoscheme reflection group $W_3/N_{\mathrm{gen}}=W_3$ is finitely
presented.
\end{enumerate}
\end{corollary}

\begin{proof}
(i) $N_{\mathrm{gen}}=\bigcap_a\ker\rho_a\subseteq\ker\rho_{(1,2,11)}=1$ by
Theorem~\ref{thm:faithful}. For the co-null set, fix $w\ne1$ in $W_3$. The entries of
$\rho_a(w)$ are rational functions of $a$ whose denominators do not vanish on the positive
orthant, so $Z_w=\{a\in\R_{>0}^3:\rho_a(w)=1\}$ is closed and algebraic there, and proper
because $(1,2,11)\notin Z_w$. A proper algebraic subset is null and $W_3$ is countable, so
$\bigcup_{w\ne1}Z_w$ is null; off it every $\rho_a$ is injective and $u_d(a)=[t^d]W_3(t)$ for
every $d$. (ii) The image
$\pi(\rho_a(W_3))\le O(3)$ is isomorphic to $W_3$ by Theorem~\ref{thm:faithful}. (iii) $W_3$ is a
right-angled Coxeter group, so it is finitely presented.
\end{proof}

\begin{remark}[a second route to Theorem~\textup{\ref{thm:masterCfalse}}]\label{rem:secondroute}
Lemma~\ref{lem:W3struct} is tuple independent, so it yields a proof of non-faithfulness that does
not use Proposition~\ref{prop:transtrick}. Since $W_3=\langle A,C\rangle\rtimes V$ with
$\langle A,C\rangle$ free on $\{A,C\}$, and $V$ always injects, the point group is injective if
and only if the images of $A$ and $C$ generate a free group on those two elements. At $(1,3,2)$
and at $(3,1,5)$ a search over reduced words in $A^{\pm1},C^{\pm1}$ in exact rationals finds the
first collision at word length $7$, and, writing $c$ for $C^{-1}$ and $a$ for $A^{-1}$,
\[
 AAcAccAAACACCA \quad\text{at }(1,3,2),\qquad AAcaaCAAACaacA \quad\text{at }(3,1,5)
\]
are reduced words of length $14$ whose images are the identity while they are nontrivial in $W_3$.
\end{remark}

\begin{remark}[the picture in dimension three]\label{rem:n3picture}
Theorems~\ref{thm:masterCfalse} and \ref{thm:faithful} are complementary and together settle the
three-dimensional case. Faithfulness holds generically, and fails on a nonempty set of rational
Class C tuples. So Class C is not the right condition: it neither implies faithfulness, by
Theorem~\ref{thm:masterCfalse}, nor is needed for it. The exceptional set is where the point
group collapses, and by Corollary~\ref{cor:necessary} that is exactly where $\rho_a$ can fail.

Nothing here decides $n\ge4$. The argument used that $W_3$ is virtually free, and that already
fails at $n=4$, though not for the reason one might expect. A triangle in $\Gamma_n$ is no
obstruction: it contributes a \emph{finite} parabolic $(\Z/2)^3$, compatible with virtual
freeness. The obstruction is an induced four-cycle. In $\Gamma_4$ the set $\{0,3,1,4\}$ induces
the cycle $0-3-1-4-0$, its non-edges being $\{0,1\}$ and $\{3,4\}$; a four-cycle is the join of
two non-adjacent pairs, so the corresponding parabolic is $D_\infty\times D_\infty$, which
contains $\Z^2$. Hence $W_4$ is not virtually free and Lemma~\ref{lem:W3struct} has no analogue
there.
\end{remark}

\begin{remark}[what the reduction does not buy]\label{rem:notbuys}
It does not decide Problem~\ref{prob:embed} for any $n\ge4$, the case $n=3$ being settled
affirmatively by Corollary~\ref{cor:three}, and by itself it does not touch
Problem~\ref{prob:alternative}: an abstract embedding of $W_n$ into $O(n)$ need not be by
reflections, and the faithfulness of the particular map $\rho_a$ is a separate question. The two
standard invariants remain undecisive. Finite subgroups do not separate: every finite subgroup of
$W_n$ is conjugate into an elementary abelian $2$-group on a clique of $\Gamma_n$, of rank at
most $\lceil(n+1)/2\rceil\le n$, while $O(n)$ contains $(\Z/2)^n$. Free subgroups do not
separate either, both groups containing free subgroups of rank two for $n\ge3$.
\end{remark}

\section{The planar case}\label{sec:plane}

Let $F=C_2\ast C_2\ast C_2$ and $\rho_\tau\colon F\to G_\tau$ the generic triangle group.

\begin{theorem}[Census; PROVED, computational input VERIFIED]\label{thm:census}
There are $33$ pairs of reduced words of length $11$, distinct in $F$, with
$\rho_\tau(w)=\rho_\tau(w')$ for every $\tau$. Outside a proper Zariski-closed subset of shape
space these are the only coincidences in the ball of radius $11$, and there are $132$ further
pairs at radius $12$. The generic sphere sizes are
\[
 1,3,6,12,24,48,96,192,384,768,1536,3039,6012,11925,23570 .
\]
\end{theorem}

The $33$ identities are polynomial identities in $\Z[p,q]$; the lower bounds come from counts in
a finite field chosen to contain the needed roots of unity and $\sqrt{-1}$, and match the upper
bounds from the identities, so the counts are exact.

\begin{theorem}[Right-triangle structure; PROVED]\label{thm:right}
For $W=\langle R_0,R_1,R_2\mid R_i^2=(R_0R_1)^2=1\rangle$ the rotation subgroup is of index $4$
and isomorphic to $\Z\wr\Z$, with translation lattice $\mathcal T=2(t-1)\Z[t^{\pm1}]$ obtained
from Crowell's exact sequence and the Bass--Serre tree of the free product. The universal
sequence is shape-independent through depth $32$, and the shape $(1,2)$ deviates at depth $33$.
If $c_T=N(a-ib)>2d$ then no deviation occurs at depth $d$.
\end{theorem}

\begin{remark}[status of the metric]\label{rem:metric}
Word length equals relaxed length plus twice the connectivity defect, $\ell_T=\ell_R+2c$. The
upper bound is PROVED. The lower bound is open and is stated as a conjecture; the closed form
$c_{\mathrm{pred}}$ computed by the explicit transducer agrees with the true defect with $0$
exceptions on the $275{,}823$ elements of the ball of radius $24$, held out at depths
$19,21,23,24$, which is HEURISTIC support and not a proof. The junction site cost is
$\max(|d_L-1|,|d_R|)$, which depends on the sign of $d_L$ and not on its magnitude alone; the
form $\max(|d_L|-1,|d_R|)$ is false on every cell with $d_L\le-2$ and $|d_R|\le|d_L|$. This is a
term inside the site-cost chain of the strand model, not an additive correction to the length
formula.
\end{remark}

\section{The growth rate $\beta_2$}\label{sec:beta}

With $q=x^2$ and $\Sigma_1$ the odd $k$-recursion sum, $1-\Sigma_1(q)$ is a Hahn--Exton
$q$-cosine, and $q^\ast=0.449453630558948046\ldots$ is its smallest positive zero.

\begin{theorem}[PROVED]\label{thm:beta}
The growth rate of the planar right-triangle group is
$\beta_2=1/\sqrt{q^\ast}=1.4916177871\ldots$, and $\beta_2<3/2$.
\end{theorem}

The proof squeezes the true count between the pure-travel and relaxed counts, which share the
exponential rate, and locates $q^\ast$ by a contour argument with exact rational winding number.
Irrationality of $\beta_2$ is open (Remark~\ref{rem:notclaimed}).

\section{Transcendence: the unconditional part}\label{sec:trans}

\begin{theorem}[PROVED]\label{thm:pc}
$\Sigma_1$ and $S_1$ are transcendental over $\Q(q)$ and have $|q|=1$ as a natural boundary.
\end{theorem}

\begin{theorem}[Travel resolvent; PROVED]\label{thm:travel}
$\Sigma_0/(1-\Sigma_1)$ is transcendental over $\Q(x)$, unconditionally.
\end{theorem}

Transcendence of $V$ is conditional on $(M)$ and $(R\!-\!J)$; transcendence of $U$ is conditional
on $(M)$, $(R\!-\!J)$ and $(T)$. Neither is claimed here.

\section{The \texorpdfstring{$q$}{q}-cosine increment}\label{sec:qcos}

\begin{theorem}[PROVED]\label{thm:qcos}
For every algebraic $q\in(0,1)$ the difference Galois group of
$G(q,z)={}_1\phi_1(0;q;q^2,z)$ over $\C(z)$ contains $SL_2$, so $G$ is hypertranscendental. The
connection constant is $C_1=1/(q;q)_\infty$, whence
$\prod_{k\ge1}z_k(q)q^{2(k-1)}=(q;q^2)_\infty$.
\end{theorem}

The general hypertranscendence statement, for all $q$ not a root of unity and over
$\bigcup_j\C(z^{1/j})$, is due to Adamczewski, Dreyfus and Hardouin and is stronger. The increment
is the removal of a transcendence hypothesis on $q$ in the $SL_2$ statement.

\section{The five directions}\label{sec:open}

\begin{problem}[Prove $(R\!-\!J)$]\label{prob:rj}
Show $\Pi_y(q_m)\ne0$ for infinitely many travel poles. This would make $V$, and with $(T)$ also
$U$, transcendental over $\Q(x)$ unconditionally, modulo $(M)$. What exists is an upper bound
$|P_{12}|w^3<2$; what is needed is a lower bound $|P_{12}(q_m)|\ge c\,w_m^{-3}$.
\end{problem}

\begin{remark}[the arithmetic structure at $(1,2,3)$]\label{rem:pingpong}
For the generic form of Problem~\ref{prob:alternative} a single faithful tuple suffices, and
by Corollary~\ref{cor:necessary} it is enough to find one at which the \emph{point group} is
faithful. The tuple $(1,2,3)$, where the point group agrees with the envelope through depth $22$,
carries an arithmetic structure that the others do not, and we record it because it is what any
attempt should use.

The two denominators are $|m_1|^2=a_1^2+a_2^2=5$ and $|m_2|^2=a_2^2+a_3^2=13$, both prime. So
$R_0,R_2,R_3$ lie in $GL_3(\Z_5)$ while $v_5(R_1)=-1$, and $R_0,R_1,R_3$ lie in $GL_3(\Z_{13})$
while $v_{13}(R_2)=-1$. The generators are orthogonal for the standard form, and
$x^2+y^2+z^2$ is isotropic over $\Q_5$ and over $\Q_{13}$, so $SO_3$ is split at both places and
its Bruhat--Tits building there is a \emph{tree} rather than a rank-two building. Finally the
graph splits: with $\Gamma_1=\{0,2,3\}$ and $\Gamma_2=\{0,1,3\}$ one has
$\Gamma_1\cap\Gamma_2=\{0,3\}$, and $1$ and $2$ are non-adjacent in $\Gamma_3$, so
\[
 W_3\;=\;(\Z/2\times D_\infty)\ \ast_{\ \Z/2\times\Z/2}\ (D_\infty\times\Z/2),
\]
the first factor $\langle R_0,R_2,R_3\rangle$ being $5$-integral and the second
$\langle R_0,R_1,R_3\rangle$ being $13$-integral, with the amalgamated
$\langle R_0,R_3\rangle$ integral at both. Both factors are faithfully represented: each contains
a $D_\infty$ acting as rotations by an angle with rational cosine $3/5$ respectively $-5/13$,
which has infinite order by Niven's theorem, and in each the extra involution reverses an axis
that the $D_\infty$ fixes. What is missing is the step from these data to injectivity of the
amalgam.

The obstacle is that this is a ping-pong on \emph{two} trees, one at each place, and a group
acting on a product of trees is outside Bass--Serre theory. That can be removed. Solving
$a_1^2+a_2^2=5^j$ and $a_2^2+a_3^2=5^k$ in distinct positive integers gives $(1,2,11)$, where
$|m_1|^2=5$ and $|m_2|^2=125$, so the whole point group lies in $O_3(\Z[1/5])$ and acts on the
\emph{single} tree at $p=5$, with $R_0,R_3$ integral and $v_5(R_1)=-1$, $v_5(R_2)=-3$. The point
group there agrees with the envelope through depth $21$ at both primes. So $(1,2,11)$, not
$(1,2,3)$, is the tuple at which the argument should be attempted: one tree, and Bass--Serre
theory applies to it. We have not carried the argument out, and record only that the arithmetic
preconditions hold.
\end{remark}

\begin{problem}[The alternative in dimension $n\ge4$]\label{prob:alternative}
Is $N_{\mathrm{gen}}=\bigcap_a\ker\rho_a$ trivial, equivalently is $\rho_a$ injective for
generic $a$? A single faithful tuple answers this affirmatively. The arithmetic form is FALSE
(Theorem~\ref{thm:masterCfalse}), which does not decide the generic form.
\end{problem}

\begin{remark}[Class C faithfulness]\label{conj:masterC}
The conjecture that $\rho_a$ is injective for every $a\in\Q_{>0}^n$ with $e(a)=t(a)=0$ and
$n\ge3$ is FALSE, by Theorem~\ref{thm:masterCfalse}. Only the generic form (a) survives.
\end{remark}

\begin{problem}[Ambient embedding]\label{prob:embed}
For $n\ge3$, does $W_n$ embed in $O(n)$? By Proposition~\ref{prop:reduce} this is equivalent to
the ambient embedding question for $\Isom(\R^n)$, and by Remark~\ref{rem:buys} it asks whether
$W_n$ has a faithful orthogonal representation of degree $n$.
\end{problem}

\begin{problem}[Finite presentation]\label{prob:fp}
For $n\ge3$, is $W_n/N_{\mathrm{gen}}$ finitely presented? A negative answer refutes
Problem~\ref{prob:alternative}(a). At $n=2$ the answer is no.
\end{problem}

\begin{problem}[Transcendence of the zeros]\label{prob:zeros}
Are the zeros $z_k(q)$ of the Hahn--Exton $q$-cosine transcendental at algebraic $q\in(0,1)$, and
is $\beta_2$ transcendental? This is the $q$-analogue of Siegel's theorem on the zeros of Bessel
functions. Even irrationality of $\beta_2$ is open.
\end{problem}

\end{document}
